% Figure: the AC-to-DC resistance ratio of a round copper wire against
% frequency, showing the sqrt(f) rise once the skin depth falls below the
% wire radius. Below the transition the current fills the wire and the
% ratio is one; above it the current crowds into a surface annulus of
% thickness delta and the effective area shrinks, so R_ac/R_dc rises. We
% plot the high-frequency approximation R_ac/R_dc = a/(2*delta) clamped to
% a floor of 1 to stitch the two regimes.
\begin{figure}[htbp]
\centering
\begin{tikzpicture}
\begin{axis}[
  width=12cm, height=6.5cm,
  xlabel={frequency (Hz, log scale)},
  ylabel={$R_{\mathrm{ac}}/R_{\mathrm{dc}}$},
  xmode=log,
  xmin=1000, xmax=1e9, ymin=0, ymax=60,
  axis lines=left,
  legend pos=north west,
  legend cell align=left,
  domain=1000:1e9, samples=300,
]
  % Copper: delta = 66.1/sqrt(f) mm approx (delta in mm, f in Hz):
  % delta[m] = 1/sqrt(pi*f*mu0*sigma), sigma=5.96e7, mu0=4e-7*pi.
  % delta[mm] = 1000/sqrt(pi*f*(4*pi*1e-7)*5.96e7) = 65.2/sqrt(f).
  % Wire radius a = 1 mm. High-f ratio ~ a/(2 delta) = 1/(2*0.0652/sqrt(f))
  %   = sqrt(f)/(0.1304). Clamp with a smooth max against 1.
  % Use ratio = sqrt( 1 + (a/(2 delta))^2 ) as a stitching form.
  \addplot[engblue, very thick] {
    sqrt( 1 + ( 1/(2*65.2/sqrt(x)) )^2 )
  };
  \addlegendentry{$1\,\mathrm{mm}$-radius wire}
  % thinner wire a = 0.3 mm, transition pushed higher.
  \addplot[engred, thick, dashed] {
    sqrt( 1 + ( 0.3/(2*65.2/sqrt(x)) )^2 )
  };
  \addlegendentry{$0.3\,\mathrm{mm}$-radius wire}
  % DC floor line
  \addplot[enggray, thin, dotted] { 1 };
  \addlegendentry{DC value}
\end{axis}
\end{tikzpicture}
\caption[AC-to-DC resistance ratio of a round wire against
frequency]{Ratio of alternating-current to direct-current resistance
for round copper wires of two radii. At low frequency the skin depth
exceeds the wire radius, the current fills the cross-section, and the
ratio holds near one. Once the skin depth falls below the radius the
current crowds into a surface layer of thickness set by the skin depth,
the effective conducting area shrinks, and the ratio climbs as the
square root of frequency. A thinner wire keeps its resistance ratio near
one to a higher frequency, the reason litz wire subdivides a conductor
into many fine strands.}
\label{fig:vol04ch10:skin-effect-resistance}
\end{figure}
